% ===================================================================
%  तालिका सं. 8 — फ़सल। — शिकार, अंगूर की फ़सल, मछली पकड़ना।
%  Traduction 2026 d'apres texto/io, controlee sur texto/fr et
%  texto/en. Consigne : voir texto/hi/10-talika-01.tex.
%
%  L'appel de note est dans le TITRE de la scene — « फ़सल (*) » — et
%  non dans un alinea : c'est ainsi que l'ido le compose.
% ===================================================================

%%K t08-noto-1 noto
\VUnotes{143.2mm}{%
\textit{(*) क्योंकि यह शब्द अस्पष्ट है : सन की फ़सल? अंगूर की? सेब की? शायद
बेहतर हो कि «} \VUgras{mesono} \textit{» काम में लाया जाए, जो} Mondo \textit{XIV,
पृ. 242 में सुझाया गया है। पर अब तक उस नई धातु को अकादमी ने अपनाया नहीं और वह
इदो आंदोलन के सामान्य नियमों के अनुसार चर्चा की प्रतीक्षा में है।}%
}

%%K t08-apar-1 apar
\VUsaut{5.0mm}
\VUtitre{15.0pt}{60}{तालिका सं. 8}
\VUsaut{3.0mm}
\VUcentre{13.2pt}{90}{\textit{पहला दृश्य।}}
\VUsaut{3.0mm}
\VUpk{10.2pt}{40}{फ़सल (*)।}
\VUsaut{4.0mm}
\VUpk{10.2pt}{40}{देहात का रूप।}
\VUsaut{3.0mm}

%%K t08-c1-01-1 p
1. --- \VUgras{गरमी} के साथ देहात ने फिर एक बार अपना रूप बदल लिया है। कूँड़ को
सौंपा गया बीज फूट पड़ा; ज़मीन से एक नन्हा हरा पौधा निकला और उसने अपनी बाली दी।
दाना बना, फिर पका, और अब फ़सल काटने के सिवा कुछ बाकी नहीं।

%%K t08-c1-02-1 p
2. --- \VUgras{जंगली फूल} खिल आए हैं। \VUgras{कॉर्नफ़्लावर} \textsuperscript{(1)},
\VUgras{पोस्ते के फूल} \textsuperscript{(2)} और \VUgras{डिजिटैलिस}
\textsuperscript{(3)} गेहूँ के खेतों के एकसार रंग पर अपने ताज़े \VUgras{फूलों}
का प्रसन्न विरोध रखते हैं।

%%K t08-c1-03-1 p
3. --- फलदार पेड़ों ने पत्ते निकाल लिए हैं; फल पक गए हैं। नन्हा
\VUgras{चेरी का पेड़} \textsuperscript{(4)} सुंदर \VUgras{चेरियों}
\textsuperscript{(5)} से लदा है, जिन्हें बच्चे तोड़ते और चाव से खाते हैं।

%%K t08-c1-04-1 p
4. --- रेवड़ अब सारा दिन बाहर बिता सकता है।

%%K t08-c1-04-2 p
\VUgras{मेमने} \textsuperscript{(6)} बढ़कर सुंदर \VUgras{भेड़ें}
\textsuperscript{(7)} और सुंदर \VUgras{मेढ़े} \textsuperscript{(8)} बन गए हैं,
घनी ऊन वाले। वे \VUgras{चरागाह} \textsuperscript{(9)} की \VUgras{घास} शांति से
चरते हैं, जिस पर \VUgras{गुलबहार} बिखरे हैं।

%%K t08-c1-04-3 p
जल्दी ही इन जानवरों की \VUgras{ऊन} उतारी जाएगी, जिससे हमारे कपड़ों का कपड़ा बनता
है।

%%K t08-tit-2 sub
\VUsaut{5.0mm}
\VUpk{10.2pt}{40}{चरवाहा।}
\VUsaut{3.4mm}

%%K t08-c1-05-1 p
5. --- \VUgras{चरवाहा} \textsuperscript{(10)} उन पर अधिक ध्यान देता नहीं लगता।
उसकी एक ही चिंता है क्षितिज पर जाता तूफ़ान, और \VUgras{गड़रिया}
\textsuperscript{(13)}, जो अपनी \VUgras{सुराही} \textsuperscript{(14)} होंठों तक
उठाए है, उससे भी अधिक बेफ़िक्र लगता है।

%%K t08-c1-06-1 p
6. --- गरमी दमघोंटू है; क्योंकि \VUgras{कुत्ता} \textsuperscript{(15)} भी ठंडा
होने आता है; वह \VUgras{नाले} \textsuperscript{(16)} का ताज़ा पानी चाटता है और
एक बेचारे नन्हे \VUgras{मेंढक} \textsuperscript{(17)} को चौंका देता है जो किनारे
की घास में दुबका था। वह उस साफ़ पानी में कूद पड़ता है जहाँ
\VUgras{कुमुदिनियाँ} \textsuperscript{(18)} खिली हैं।

%%K t08-c1-07-1 p
7. --- एक \VUgras{खंजन} \textsuperscript{(19)} उस नन्हे \VUgras{सोते}
\textsuperscript{(20)} के पास नहा रहा है जो दो चट्टानों के बीच से फूटता है और
\VUgras{काँटेदार झाड़ियों} \textsuperscript{(21)} तथा \VUgras{नागफनी की झाड़ियों}
\textsuperscript{(22)} के नीचे छिप जाता है।

%%K t08-tit-3 sub
\VUsaut{5.0mm}
\VUpk{10.2pt}{40}{फ़सल की कटाई।}
\VUsaut{3.4mm}

%%K t08-c1-08-1 p
8. --- देहात के बच्चों के लिए गेहूँ की कटाई बड़े मनोरंजन का समय है। वे प्रसन्न
\VUgras{कलाबाज़ियाँ} \textsuperscript{(24)} खाते हैं \VUgras{पुआल के ढेरों}
\textsuperscript{(23)} पर, \VUgras{काटने वालों} \textsuperscript{(26)} की मदद
करते हैं, और जब वे \VUgras{गेहूँ का खेत} \textsuperscript{(27)} काटते हैं अपनी
\VUgras{हँसियों} \textsuperscript{(25)} से, जो \VUgras{सान}
\textsuperscript{(30)} पर ख़ूब तेज़ की गई हैं, तब बच्चे गेहूँ बटोरकर
\VUgras{पूले} \textsuperscript{(31)} या \VUgras{छोटी गड्डियाँ} \textsuperscript{(32)}
बाँधते हैं।

%%K t08-c1-09-1 p
9. --- कभी-कभी उनमें से बड़ों को \VUgras{दराँती} \textsuperscript{(12)} से वे
\VUgras{सुनहरे डंठल} \textsuperscript{(28)} काटने दिए जाते हैं जो दानों से भरी
भारी \VUgras{बालियाँ} \textsuperscript{(29)} उठाए हैं; और छोटे, \VUgras{पशुओं}
\textsuperscript{(33)} की रखवाली करते हुए, जो \VUgras{मैदान}
\textsuperscript{(34)} में बड़े \VUgras{लिंडेन पेड़} \textsuperscript{(35)} की
छाँह में चरते हैं, अपनी \VUgras{जाली} \textsuperscript{(36)} में सुंदर \VUgras{तितलियाँ}
पकड़ते हैं।

%%K t08-c1-09-2 p
कुछ तो \VUgras{गाड़ी} \textsuperscript{(37)} पर चढ़कर वे पूले जमाते हैं जो उन्हें
\VUgras{पंजे} \textsuperscript{(38)} पर उठाकर दिए जाते हैं।

%%K t08-c1-10-1 p
10. --- सारे \VUgras{मैदान} \textsuperscript{(39)} में, जहाँ से \VUgras{चंडूलें}
\textsuperscript{(41)} उड़ती हैं, बड़ी चहल-पहल है। सब फ़सल में लगे हैं। पर जहाँ
ग़रीब पुराने औज़ार ही काम में लाते हैं — \VUgras{हँसियाँ, दराँतियाँ} और
\VUgras{मुँगरी} \textsuperscript{(40)} —, वहाँ धनी ज़मींदार अपना गेहूँ या अपनी
\VUgras{राई} \textsuperscript{(43)} \VUgras{काटने की मशीन} \textsuperscript{(42)} से
कटवाते हैं। फिर, पूलों को खलिहान तक ढोए बिना, वहीं बड़े \VUgras{ढेर}
\textsuperscript{(44)} लगा दिए जाते हैं।

%%K t08-c1-11-1 p
11. --- फिर एक \VUgras{गहाई की मशीन} \textsuperscript{(47)} लाई जाती है, जिसे एक
\VUgras{चलता इंजन} \textsuperscript{(45)} एक \VUgras{पट्टे} \textsuperscript{(46)}
के ज़रिए चलाता है, और इस तरह दाना \VUgras{पुआल} से अलग हो जाता है, उसी खेत में
जहाँ वह बोया गया था।

%%K t08-tit-4 sub
\VUsaut{5.0mm}
\VUpk{10.2pt}{40}{तूफ़ान।}
\VUsaut{3.4mm}

%%K t08-c1-12-1 p
12. --- फ़सल के समय प्रायः प्रचंड \VUgras{तूफ़ान} \textsuperscript{(53)} टूट पड़ते
हैं। पहले दूर से \VUgras{बादलों की गड़गड़ाहट} सुनाई देती है; एक तेज़
\VUgras{पछुआ हवा} \textsuperscript{(54)} उठती है और \VUgras{जंगल}
\textsuperscript{(49)} के सबसे बड़े पेड़ों को कराहते हुए झुका देती है। काले
\VUgras{बादल} \textsuperscript{(56)} आकाश पर चढ़ आते हैं; उन्हें
\VUgras{बिजली} \textsuperscript{(55)} की चौंधियाती टेढ़ी रेखाएँ चीरती हैं।
\VUgras{बारिश} और कभी \VUgras{ओले} गिरने लगते हैं, और कटने को तैयार फ़सल को
बिछा देते हैं।

%%K t08-c1-13-1 p
13. --- प्रायः बिजली किसी इमारत में आग लगा देती है। पिछले साल, उदाहरण के लिए,
\VUgras{पवनचक्की} \textsuperscript{(50)} की छत जलकर राख हो गई और उसके दो
\VUgras{पंख} \textsuperscript{(51)} चूर-चूर हो गए।

%%K t08-c1-14-1 p
14. --- कुछ घंटों बाद शांति लौट आती है। क्षितिज पर एक चमकीला \VUgras{इंद्रधनुष}
\textsuperscript{(52)} दिखाई देता है, और प्रकृति वह जीवन फिर उठा लेती है जो थोड़ी
देर के लिए रुक गया था। जो देहाती \VUgras{झोपड़ियों} \textsuperscript{(48)} में
आ छिपे थे, वे काम पर लौटते हैं और अभी-अभी हुए नुक़सान की मरम्मत में हिम्मत से
जुट जाते हैं।

%%K t08-tit-5 sub
\VUsaut{6.0mm}
\VUcentre{13.2pt}{90}{\textit{दूसरा दृश्य।}}
\VUsaut{3.6mm}

%%K t08-tit-6 sub
\VUpk{10.2pt}{40}{शिकार। --- अंगूर की फ़सल।}
\VUsaut{1.4mm}
\VUpk{10.2pt}{40}{मछली पकड़ना।}
\VUsaut{4.0mm}

%%K t08-c2-01-1 p
1. --- \VUgras{अगस्त} के अंत या \VUgras{सितंबर} के मध्य में, इलाक़े के अनुसार,
शिकार का मौसम खुलता है। सूरज की पहली किरणें \VUgras{छिपकलियों}
\textsuperscript{(2)} को उनके बिलों से बाहर लाती भी नहीं कि \VUgras{शिकारी}
\textsuperscript{(5)} अपने \VUgras{कुत्तों} \textsuperscript{(4)} के साथ निकल
पड़ता है।

%%K t08-c2-02-1 p
2. --- उसने मोटे कपड़े का अपना मज़बूत \VUgras{शिकारी सूट} \textsuperscript{(9)}
पहना है, और चमड़े के \VUgras{पट्टे}, जो उसे उस गीली घास में चलने देंगे जहाँ
\VUgras{भेक} \textsuperscript{(1)} छिपे रहते हैं; उसने अपनी \VUgras{बंदूक}
\textsuperscript{(6)} और अपना \VUgras{शिकार-थैला} \textsuperscript{(10)} लिया है,
और चल दिया है।

%%K t08-c2-03-1 p
3. --- शिकार की उमंग उसे एक अंगूर-बाग़ के बीच ले आती है जहाँ अंगूर की फ़सल पूरे
ज़ोर पर है। अंगूर उगाने वाले के बच्चे, जो आम तौर पर सड़क के किनारे बकरियाँ
चराते हैं, आज उन्हें जहाँ जी चाहे चरने दे रहे हैं, और एक बूढ़े \VUgras{बकरे}
\textsuperscript{(3)} को खूँटे से बाँधकर — जो अपनी आज़ादी का लाभ उठाकर निश्चय ही
भाग जाता — वे भी अपने हिस्से का मज़ा ले रहे हैं। एक \VUgras{अंगूर की टोकरी}
\textsuperscript{(12)} से एक गुच्छा लेता है और उसका \VUgras{रस} \textsuperscript{(14)}
एक \VUgras{प्याले} \textsuperscript{(13)} में निचोड़कर मस्त (बिना ख़मीर की शराब)
बनाने में मन बहलाता है। दूसरा अपने सिर पर \VUgras{पत्तों का मुकुट}
\textsuperscript{(15)} रख रहा है जो उसने अभी गूँथा है। उनके पास ढीठ
\VUgras{गौरैयाँ} \textsuperscript{(16)} अंगूर चुराने कोल्हू तक चली आती हैं।

%%K t08-c2-04-1 p
4. --- जब तक बच्चे मन बहलाते हैं, आदमी काम करते हैं।
\VUgras{अंगूर तोड़ने वाले} \textsuperscript{(18)} अंगूर \VUgras{पीठ की टोकरियों}
\textsuperscript{(17)} में तहख़ाने तक ले जाते हैं; एक \VUgras{पीपासाज़}
\textsuperscript{(19)} \VUgras{पीपों} \textsuperscript{(20)} की मरम्मत करता है,
देखता है कि \VUgras{फट्टे} \textsuperscript{(22)} और \VUgras{ढक्कन}
\textsuperscript{(21)} ठीक हैं, और टूटे \VUgras{छल्ले} \textsuperscript{(23)}
बदलता है। वे बड़े \VUgras{कुंड} \textsuperscript{(25)} भी उसी ने बनाए थे, जिनमें
ख़मीर उठाने को \VUgras{अंगूर} \textsuperscript{(26)} डाले जाते हैं।

%%K t08-c2-05-1 p
5. --- ख़मीर उठ चुकने पर शराब निकाल ली जाती है और बचा हुआ \VUgras{कोल्हू}
\textsuperscript{(28)} पर रखा जाता है; बड़े \VUgras{पेंच} \textsuperscript{(29)} को
धीरे-धीरे कसकर रस निचोड़ा जाता है, जो एक \VUgras{टब} \textsuperscript{(32)} में
बहता है और और भी \VUgras{शराब} \textsuperscript{(31)} देता है।

%%K t08-tit-8 sub
\VUsaut{6.0mm}
\VUpk{10.2pt}{40}{बीयर और सिडर।}
\VUsaut{3.4mm}

%%K t08-c2-06-1 p
6. --- \VUgras{तहख़ाने} \textsuperscript{(27)} की दीवार पर \VUgras{हॉप्स}
\textsuperscript{(30)} की एक शाख चढ़ती है। यहाँ वह केवल सजावटी पौधा है, पर कुछ
देशों में, जहाँ बेल बहुत कम होती है, \VUgras{बीयर} बनाने के लिए हॉप्स उगाए जाते
हैं।

%%K t08-c2-07-1 p
7. --- नन्हा \VUgras{सेब का पेड़} \textsuperscript{(33)} सेबों से लदा है, जिनसे
पकने पर बढ़िया \VUgras{सिडर} बनेगा। अपने \VUgras{पिंजरे} \textsuperscript{(34)}
में \VUgras{कनारी} \textsuperscript{(35)} और \VUgras{गोल्डफ़िंच}
\textsuperscript{(36)} उन गौरैयों को ईर्ष्या से देखते हैं जो आज़ाद उछल-कूद कर रही
हैं।

%%K t08-c2-08-1 p
8. --- जहाँ तक नज़र जाती है, पहाड़ियों की ढलान पर, \VUgras{बेल के खूँटों}
\textsuperscript{(40)} की क़तारें खड़ी हैं, जो शानदार \VUgras{बेल के ठूँठों}
\textsuperscript{(39)} को सँभाले हैं, \VUgras{गुच्छों} से लदे हुए।

%%K t08-c2-08-2 p
सारा साल \VUgras{अंगूर उगाने वाले} \textsuperscript{(42)} ने अपने
\VUgras{अंगूर-बाग़} \textsuperscript{(38)} की देखभाल में मेहनत की है, और अब उसे
अच्छी फ़सल से उसका पुरस्कार मिलता है। \VUgras{तोड़ने वालियाँ}
\textsuperscript{(41)} उन कुंडों को भरती हैं जो \VUgras{ख़च्चरों}
\textsuperscript{(43)} से खींची गाड़ी पर रखे हैं, और सब प्रसन्न हैं।

%%K t08-tit-9 sub
\VUsaut{5.0mm}
\VUpk{10.2pt}{40}{मछली पकड़ना।}
\VUsaut{3.4mm}

%%K t08-c2-09-1 p
9. --- यही बात \VUgras{बंसीवालों} \textsuperscript{(44)} की है, जो पौ फटते ही
अपनी \VUgras{बंसियों} \textsuperscript{(45)} के साथ पानी के किनारे आ जमे हैं।

%%K t08-c2-09-2 p
\VUgras{मछलियाँ} \textsuperscript{(47)} काँटा पकड़ लेती हैं ज्यों ही वे
अपनी \VUgras{डोरी} \textsuperscript{(46)} पानी में डालते हैं, और उन्हें
\VUgras{चारा} बदलने का समय भी नहीं मिलता कि डोरी के सिरे पर एक नया शिकार तड़प
रहा होता है।

%%K t08-c2-09-3 p
फिर भी \VUgras{जाल वाला मछुआ} \textsuperscript{(48)}, जो अपनी \VUgras{नाव}
\textsuperscript{(49)} से \VUgras{नदी} \textsuperscript{(50)} के बीच जाल फेंकता
है, कहीं अधिक पकड़ता है।

%%K t08-c2-10-1 p
10. --- शरद् अच्छी सैर का मौसम है : पैदल, \VUgras{घोड़े} \textsuperscript{(52)}
पर, \VUgras{साइकिल} \textsuperscript{(53)} से, \VUgras{मोटर} \textsuperscript{(54)}
से। \VUgras{सड़कें} \textsuperscript{(51)} साफ़ हैं और लोग आख़िरी अच्छे दिनों का
भरपूर लाभ उठाते हैं, क्योंकि \VUgras{अबाबीलें} \textsuperscript{(56)} पहले ही छतों
पर जुटने लगी हैं, गरम देशों को उड़ जाने के लिए। बूढ़े \VUgras{अख़रोट के पेड़}
\textsuperscript{(57)} के पत्ते पीले पड़ रहे हैं, \VUgras{अख़रोट} गिर रहे हैं, और एक
\VUgras{झुरमुट} \textsuperscript{(58)} में खड़े ऊँचे \VUgras{पेड़}, जो
\VUgras{ऊँची भूमि} \textsuperscript{(59)} पर है, जल्दी ही नंगे हो जाएँगे।

%%K t08-c2-11-1 p
11. --- \VUgras{सूरज} \textsuperscript{(60)} देर से उगता है, दिन छोटे होते जाते
हैं, \VUgras{सुबह} एक अच्छी-ख़ासी ठंडक महसूस होने लगती है, और
\VUgras{चाहों} \textsuperscript{(61)} के झुंड अभी से हमारी निर्मम जलवायु छोड़कर
जा रहे हैं।
\VUsaut{6.0mm}
\VUfilet{22mm}
